\documentclass[../../../main]{subfiles}

\ifSubfilesClassLoaded{%
    \zexternaldocument*{../../1.Category-Theory/exported/aux/main}%
    \zexternaldocument*{1.Abelian-categories}%
}{}

\begin{document}

\section{Complexes and cohomology}

Throughout this chapter we will be working over a fixed \hyperref[def:abelian-category]{abelian category} \(\category{A}\).

\begin{definition}[Complex] \label{def:complex}
  A sequence of objects and morphisms
  \begin{center}
    \begin{tikzcd}
      \cdots \arrow[r, "d_{i-2}"] & C_{i-1} \arrow[r, "d_{i-1}"] & C_{i} \arrow[r, "d_{i}"] & C_{i+1} \arrow[r, "d_{i+1}"] & \cdots
    \end{tikzcd}
  \end{center}
  in \(\category{A}\) is a \textbf{chain complex} if for every \(i \in \mathbb{N} \quad d_{i-1} \circ d_{i} = 0\).
  We denote a chain complex as \((C_{\bullet}, d_{\bullet})\) or simply as \(C_{\bullet}\).

  By reversing the order, we say that a sequence
  \begin{center}
    \begin{tikzcd}
      \cdots \arrow[r, "d^{i+2}"] & C^{i+1} \arrow[r, "d^{i+1}"] & C^{i} \arrow[r, "d^{i}"] & C^{i-1} \arrow[r, "d^{i-1}"] & \cdots
    \end{tikzcd}
  \end{center}
  is a \textbf{cochain complex} if for every \(i \in \mathbb{N} \quad d^{i} \circ d^{i+1} = 0\).
  We denote a cochain complex as \((C^{\bullet}, d^{\bullet})\) or simply as \(C^{\bullet}\).
\end{definition}

\begin{remark}
  In purelly algebraic terms, chains are cochains are \textit{equivalent}.
  We will adopt the cochain notation, but all proofs and definitions translate to the chain notation.
\end{remark}


\begin{definition}[Exact complex] \label{def:exact-complex}
  A \hyperref[def:complex]{complex}
  \begin{center}
    \begin{tikzcd}
      \cdots \arrow[r, "d^{i+2}"] & C^{i+1} \arrow[r, "d^{i+1}"] & C^{i} \arrow[r, "d^{i}"] & C^{i-1} \arrow[r, "d^{i-1}"] & \cdots
    \end{tikzcd}
  \end{center}
  in \(\category{A}\) is \textbf{exact at \(C^{i}\)} (or simply at position \(i\)) if
  \begin{equation*}
    \coker d^{i+1} \circ \ker d^{i} = 0
  \end{equation*}

  A complex is \textbf{exact} if it is exact at every position.
\end{definition}

\begin{definition}[Short exact sequence] \label{def:short-exact-sequence}
  An \hyperref[def:exact-complex]{exact complex} of the form
  \begin{center}
    \begin{tikzcd}
      0 \arrow[r] & X \arrow[r] & Y \arrow[r] & Z \arrow[r] & 0
    \end{tikzcd}
  \end{center}
  is called a \textbf{short exact sequence}.

  \hyperref[def:exact-complex]{Exact complexes} of the form \(0 \to X \to Y \to Z\) (resp. \(X \to Y \to  Z \to 0\)) are called left (resp. right) short exact sequences.
\end{definition}


\begin{remark} \label{remark:exactness-mono-epi}
  For any morphism \(f: X \to Y\) in an \(\category{A}\) we get \hyperref[def:short-exact-sequence]{short exact sequences}
  \begin{center}
    \begin{tikzcd}
      0 \arrow[r] & \ker f \arrow[r] & X \arrow[r] & \im f \arrow[r] & 0 \\
      0 \arrow[r] & \im f \arrow[r] & Y \arrow[r] & \coker f \arrow[r] & 0
    \end{tikzcd}
  \end{center}

  Therefore by \autoref{lemma:partial-sequence-additive} the sequence \(0 \to X \xrightarrow{f} Y\) (resp. \(X \xrightarrow{f} Y \to 0\)) is \hyperref[def:exact-sequence]{exact} if and only if \(f\) is a \hyperref[def:mono-epi-morphism]{monomorphism} (resp. \hyperref[def:mono-epi-morphism]{epimorphism}).
\end{remark}

\end{document}
